% ============================================================
\section{การบูรณาการระบบ (System Integration)}
\label{sec:system_integration}
% ============================================================

\subsection{สถาปัตยกรรมการบูรณาการ}
\label{subsec:integration_architecture}

การบูรณาการระบบเชื่อมต่อ Subsystem ทั้งหมด 4 ส่วน
ได้แก่ Mechanical, Electrical, Firmware และ
External Application เข้าด้วยกันผ่าน Interface
ที่กำหนดไว้ใน System Architecture

% [รูปที่ XX: Integration Diagram ของระบบทั้งหมด]
% \begin{figure}[H]
%   \centering
%   \includegraphics[width=0.95\textwidth]{images/integration_diagram.png}
%   \caption{Integration Diagram ของระบบทั้งหมด}
%   \label{fig:integration_diagram}
% \end{figure}

% ------------------------------------------------------------
\subsection{การเชื่อมต่อ STM32 กับ Motor และ Encoder}
\label{subsec:stm32_motor_encoder_connection}
% ------------------------------------------------------------

STM32G474RE เชื่อมต่อกับ Motor และ Encoder
ผ่าน Peripheral ดังนี้

\begin{itemize}
  \item \textbf{Encoder}: AMT-103 ต่อกับ TIM3
        ในโหมด \texttt{TIM\_ENCODERMODE\_TI12}
        (4x Quadrature) ผ่าน Pin PB4 (CH A)
        และ PB5 (CH B) Counter Range 0--65535
        Resolution = $2048 \times 4 \times 70
        = 573{,}440$ counts/rev ที่ Output shaft

  \item \textbf{PWM}: TIM8 CH2 บน Pin PA8
        Prescaler = 169, Period = 50
        ได้ $f_{PWM} = 170\,\text{MHz}
        / (170 \times 51) \approx 19.6\,\text{kHz}$
        ทิศทางควบคุมผ่าน GPIO PC6 (DIR)

  \item \textbf{Motor Driver}: Cytron MD10C
        รับ PWM และ DIR จาก STM32 โดยตรง
        ผ่าน Optocoupler Module เพื่อแยก Ground
\end{itemize}

ผลการทดสอบการอ่านค่า Encoder และสัญญาณ PWM
แสดงดังรูปที่~\ref{fig:encoder_pwm_test}

% [รูปที่ XX: ผลการทดสอบการอ่านค่า Encoder และสัญญาณ PWM]
% \begin{figure}[H]
%   \centering
%   \includegraphics[width=0.88\textwidth]{images/encoder_pwm_test.png}
%   \caption{ผลการทดสอบการอ่านค่า Encoder และสัญญาณ PWM
%            แสดง Count และมุมที่คำนวณได้แบบ Real-time}
%   \label{fig:encoder_pwm_test}
% \end{figure}

% ------------------------------------------------------------
\subsection{การเชื่อมต่อ Joystick}
\label{subsec:joystick_connection}
% ------------------------------------------------------------

Joystick เชื่อมต่อกับ STM32 ผ่าน ESP32-DEVKIT-V1
ซึ่งทำหน้าที่เป็น Bluetooth Receiver
รับสัญญาณจาก Wireless Joystick แล้วแปลงเป็น
Command Character ส่งต่อผ่าน USART3
ที่ Baud Rate 115200\,bps, 8N1

การทดสอบครอบคลุมทุกโหมดการทำงาน
ได้แก่ Manual Mode และ Auto Mode
โดย E-Stop ทำงานได้ทุกโหมดตามข้อกำหนด J1

\begin{table}[H]
  \centering
  \caption{ผลการทดสอบ Joystick แต่ละโหมด}
  \label{tab:joystick_test}
  \begin{tabular}{lllc}
    \toprule
    โหมด & ปุ่มที่ทดสอบ & พฤติกรรมที่สังเกตได้ & ผล \\
    \midrule
    Manual
      & หมุนซ้าย (\texttt{L})
      & แขนหมุนซ้ายตาม Step ที่กำหนด
      & ผ่าน \\
    Manual
      & หมุนขวา (\texttt{R})
      & แขนหมุนขวาตาม Step ที่กำหนด
      & ผ่าน \\
    Manual
      & Gripper Pick (\texttt{S})
      & Gripper ทำ Sequence Pick ครบ
      & ผ่าน \\
    Manual
      & Gripper Place (\texttt{G})
      & Gripper ทำ Sequence Place ครบ
      & ผ่าน \\
    ทุกโหมด
      & Home (\texttt{A} Hold)
      & แขนเคลื่อนกลับตำแหน่ง 0°
      & ผ่าน \\
    ทุกโหมด
      & E-Stop (\texttt{P})
      & มอเตอร์หยุดทันที PWM = 0
      & ผ่าน \\
    ทุกโหมด
      & Reset E-Stop (\texttt{P} ซ้ำ)
      & ระบบกลับสู่ Ready State
      & ผ่าน \\
    Manual
      & สลับ Coarse/Fine (\texttt{M})
      & Jog Mode สลับระหว่าง Step/Jog
      & ผ่าน \\
    \bottomrule
  \end{tabular}
\end{table}

% ------------------------------------------------------------
\subsection{การเชื่อมต่อ Modbus กับ Base System}
\label{subsec:modbus_connection}
% ------------------------------------------------------------

STM32G474RE ทำงานเป็น Modbus RTU Slave
บน LPUART1 ที่ Slave Address 21,
Baud Rate 19200\,bps, 8E1
รองรับ Function Code 0x03 (Read Holding Registers)
และ 0x06 (Write Single Register)
ทดสอบการสื่อสารด้วย QModMaster

\paragraph{Modbus Register Map}

\begin{table}[H]
  \centering
  \caption{Modbus Register Map (Slave Address 21)}
  \label{tab:modbus_map}
  \begin{tabular}{lllp{5.5cm}}
    \toprule
    Address (Hex) & ประเภท & ชื่อ & คำอธิบาย \\
    \midrule
    \multicolumn{4}{l}{\textit{Write Registers (Master $\rightarrow$ Slave)}} \\
    \midrule
    0x00 & R    & Device ID    & ค่าคงที่ 22881 (``YA'') Heartbeat \\
    0x01 & R/W  & Command Bits &
      Bit0: Go Home,
      Bit1: Jog/Gripper Manual,
      Bit2: Auto Mode,
      Bit3: Set Home (Reset Encoder),
      Bit4: Test Mode \\
    0x02 & R/W  & Gripper Cmd  &
      0=Up, 1=Down, 2=Open, 4=Close
      (ใช้คู่กับ Bit1 ของ 0x01) \\
    0x03 & R/W  & Gripper Seq  & 1=Pick Sequence, 2=Place Sequence \\
    0x05 & R/W  & Jog Step     & ระยะ Jog (deg, signed int16) \\
    0x24 & R/W  & Target Pos   & ตำแหน่งเป้าหมาย (deg, signed int16) \\
    0x25 & R/W  & Safety Bits  & Bit0: E-Stop, Bit1: Reset E-Stop \\
    \midrule
    \multicolumn{4}{l}{\textit{Read Registers (Slave $\rightarrow$ Master)}} \\
    \midrule
    0x26 & R    & Reed Sensors & สถานะ Reed Switch (ปัจจุบัน = 0) \\
    0x27 & R    & Task Status  &
      Bit1: Position Mode Active,
      Bit2: Ghost Move Active \\
    0x28 & R    & Current Pos  & ตำแหน่งปัจจุบัน ($\times$10, signed int16) \\
    0x29 & R    & Current Speed& ความเร็วปัจจุบัน ($\times$10, signed int16) \\
    0x30 & R    & Acceleration & ความเร่งปัจจุบัน (ปัจจุบัน = 0) \\
    0x31 & R    & Safety Status& 0=Normal, 1=E-Stop Active \\
    \bottomrule
  \end{tabular}
\end{table}

หมายเหตุ: ค่าตำแหน่งที่ Register 0x28 สเกล $\times 10$
เพื่อให้ส่งค่า Decimal ได้ผ่าน Integer Register
เช่น ตำแหน่ง 90.5° ส่งเป็น 905

% ------------------------------------------------------------
\subsection{แผนการทดสอบ End-to-End}
\label{subsec:end_to_end_test_plan}
% ------------------------------------------------------------

\begin{table}[H]
  \centering
  \caption{แผนการทดสอบ End-to-End}
  \label{tab:e2e_plan}
  \begin{tabular}{clp{4.5cm}lc}
    \toprule
    \# & Test Case & เกณฑ์วัดผลสำเร็จ & Requirements & สถานะ \\
    \midrule
    1 & Encoder + PWM Verification
      & อ่านค่ามุมได้ถูกต้อง $\pm$[THRESH]°,
        PWM ควบคุมความเร็วได้
      & M4, A2
      & ผ่าน \\
    2 & E-Stop Hardware Interlock
      & กด E-Stop มอเตอร์หยุดทันที,
        STM32 ยังมีไฟ
      & E6, E7
      & ผ่าน \\
    3 & Joystick Manual Control
      & ควบคุมทุกปุ่มได้ถูกต้อง
        Response Time $<$100\,ms
      & E8, J1
      & ผ่าน \\
    4 & Modbus Read/Write
      & QModMaster อ่าน 0x28 ได้ค่าตำแหน่งจริง,
        เขียน 0x24 แล้วแขนเคลื่อนที่
      & E1
      & ผ่าน \\
    5 & PID Closed-Loop (No Load)
      & Settling Time $\leq$0.5\,s,
        Overshoot $\leq$1\%
      & P5, P6
      & รอ Tune \\
    6 & PID Closed-Loop (2\,kg Load)
      & Settling Time $\leq$0.5\,s,
        Overshoot $\leq$1\%
      & P5, P6, P7
      & รอ Tune \\
    7 & Full Auto Sequence (4 ชิ้น)
      & Cycle Time $\leq$35\,s
        ครบ 4 ชิ้นใน 1 รอบ
      & P1, P2
      & รอ Demo \\
    8 & Precision Test (100 รอบ)
      & Repeatability $\leq \pm$0.1°
      & A2
      & รอ Demo \\
    9 & Continuous Rotation 540°
      & สายไม่พันกัน ระบบทำงานปกติ
      & M9
      & รอ Demo \\
    \bottomrule
  \end{tabular}
\end{table}